\section{Introduction}\label{sec:introduction}

We prove Sogge's local smoothing conjecture for the Euclidean wave
equation in every spatial dimension \(n\ge2\). A wave arriving from a
thin spherical shell may have a much larger \(L^p\) norm at the instant
when its rays come to a focus than it had initially. Sogge asked how
many derivatives of the initial data are needed to control the norm
averaged over a time interval \cite[Section~4, Equation~(4.5)]{Sogge1991}.

Write \(U(T)=e^{-\mathrm iT|D|}\) for the half-wave propagator, with
\(|D|=\sqrt{-\Delta}\). Put \(I=[1,2]\), \(J=(1-\Delta)^{1/2}\), and
\[
 s_+(p)=\max\left\{0,\frac{n-1}{2}-\frac np\right\}.
\]

\begin{theorem}[Euclidean local smoothing]\label{thm:local-smoothing}
For \(n\ge2\), \(2\le p<\infty\) and \(\varepsilon>0\),
\begin{equation}\label{eq:sogge}
 \|U(T)f\|_{L^p(\RR^n\times I)}
     \le C_{n,p,\varepsilon}\|J^{s_+(p)+\varepsilon} f\|_{L^p(\RR^n)} .
\end{equation}
At \(p=2\) one can take \(\varepsilon=0\).
\end{theorem}

At a fixed nonzero time the sharp Sobolev order is
\((n-1)(1/2-1/p)\) \cite{Peral1980,Miyachi1980,SSS1991}.
Thus for \(p\ge\pc=2n/(n-1)\), averaging in time gains every amount
less than \(1/p\) over the fixed-time estimate. For \(2<p<\pc\),
arbitrarily small positive regularity suffices. A focusing wave shows
that \(s_+(p)\) is the smallest possible threshold
(Proposition~\ref{prop:sharpness}). We state the wave estimates for complex
Schwartz data; density gives the indicated extensions, and conjugation
gives the other sign of the half wave. Unqualified \(L^p\) norms of the
data are spatial norms.

\papercontents

\subsection{Consequences for Fourier analysis and geometry}

\paragraph{Bochner--Riesz means and Stein's square function.}
The Bochner--Riesz means \(B_R^\alpha=(1-|D|^2/R^2)_+^\alpha\)
are uniformly bounded on \(L^p\) for \(1\le p\le\infty\) and every
\(\alpha>\max\{0,n|1/p-1/2|-1/2\}\), as conjectured. For
\(2\le p<\infty\) the same orders give a strong \(L^p\) bound for
the supremum over all radii (Theorem~\ref{thm:bochner-riesz}).
Stein's square function measures their variation with radius~\cite{Stein1958Square}.
Keeping this variation inside a square integral gives its conjectured
bounds for \(2\le p<\infty\) and \(\alpha>1/2+s_+(p)\)
(Theorem~\ref{thm:stein-square}).
The connection with the wave is that averaging in time against an
oscillatory weight selects a thin shell in frequency. Sogge used this
observation to pass from local smoothing to summability of Fourier
integrals \cite[Equations~(4.5)--(4.6)]{Sogge1991}; for the reductions
at fixed and variable radii see \cite[Propositions~39 and~40]{BHS2021}.

\paragraph{Fourier restriction and Schr\"odinger waves.}
Give an incoming spherical wave an amplitude that varies with
direction. Near the focus, rescaling recovers the Fourier extension
of this amplitude from the sphere. We obtain the full interior range
of global Fourier extension estimates \(L^a\to L^b\): for
\(1<a\le\infty\), \(b<\infty\), it is given by
\[
 b>\frac{2n}{n-1},\qquad b>\frac{n+1}{n-1}\,a',
       \qquad \frac1a+\frac1{a'}=1 .
\]
Slicing gives the same bounds for the truncated cone in \(\RR^{n+1}\)
(Corollary~\ref{cor:spherical-restriction}). Rescaling a small cap
gives the compact paraboloid estimate \eqref{eq:paraboloid-extension-mixed}
and local smoothing for Schr\"odinger waves in \(n-1\) spatial dimensions,
with the exact Sobolev order above their critical exponent
(Corollary~\ref{cor:parabolic}).

\paragraph{Kakeya and Nikodym sets.}
Random signs on
disjoint frequency caps turn the square of the extension into a
bound for overlaps of tubes. This proves the Kakeya maximal
conjecture: the operator that takes \(f\) to the largest average of
\(|f|\) on a unit \(\delta\)-tube in each direction has norm
\(O_\varepsilon(\delta^{-\varepsilon})\) from \(L^n(\RR^n)\) to
\(L^n(S^{n-1})\) for every \(\varepsilon>0\). In particular, every subset of \(\RR^n\) containing
a unit line segment in each direction has Hausdorff dimension \(n\)
(Corollary~\ref{cor:quant-kakeya}). A projective change of variables
gives the corresponding Nikodym maximal estimate for tubes through a
given point (see \eqref{eq:kakeya-nikodym}).

\subsection{Development of the local smoothing problem}\label{subsec:history}

The question grew out of Bourgain's theorem on maximal averages over
circles \cite{Bourgain1986Circles}. At high frequency the Fourier
transform of a circle of radius \(T\) has two leading terms, with phases
\(T|\xi|\) and \(-T|\xi|\). This explains why a spacetime estimate for
waves can help bound the supremum of circular averages over radii.
Sogge used this connection to prove a positive smoothing gain for every
\(p>2\) in two space dimensions \cite[Theorem~4.1]{Sogge1991}. Mockenhaupt, Seeger
and Sogge gave a proof of the circular maximal theorem using angular
square functions \cite{MSS1992}, and developed local smoothing for
Fourier integral operators under a curvature condition on the
associated family of cones \cite{MSS1993}.

A persistent difficulty is to put waves moving in different directions
back together. Wolff decomposed a neighborhood of the cone into thin
angular sectors and controlled their sum by induction on scales. This
gave every gain below \(1/p\) for \(p>74\) in the plane
\cite{Wolff2000}; {\L}aba and Wolff extended the argument to higher
dimensions \cite{LabaWolff2002}. Bourgain and Demeter's cone decoupling
theorem gives this gain for \(p\ge 2(n+1)/(n-1)\) in every dimension
\cite[Theorem~1.2]{BourgainDemeter2015}. A different argument using
radial Fourier multipliers attained the exact gain \(1/p\), with no
additional Sobolev loss, for \(n\ge4\) and \(p>2(n-1)/(n-3)\)
\cite[Theorem~1.2]{HeoNazarovSeeger2011}. In the plane, Guth, Wang
and Zhang reached the critical exponent \(p=4\) by controlling the
\(L^4\) norm of a wave by the spacetime square function of its angular
pieces \cite[Theorems~1.1 and~1.2]{GWZ2020}. For the canonical
sectors of a truncated cone, Cushman, Demeter and Wu give the sharp
power \(1/4\) of the logarithm in this square function inequality
\cite[Theorem~1.4]{CushmanDemeterWu2026}. To obtain a bound in
terms of \(\|f\|_4\), one must also estimate the square function from
the initial data.

In higher dimensions the distinction between a partial gain at \(\pc\)
and a full gain above \(\pc\) is substantial. At \(n=p=3\), Gao, Liu,
Miao and Xi obtained every gain below \(2/9\) by broad Fourier
restriction estimates, compared with the conjectured \(1/3\)
\cite[Corollary~1.3]{GaoLiuMiaoXi2023}. Gan, He, Li and Wu proved the
full gain up to an arbitrarily small loss for \(p\ge10/3\) in
\(\RR^3\) and on compact Riemannian three-manifolds
\cite[Theorem~1.3]{GHLW2026}; see also Wang's study of wave packet
density \cite{Wang2026v4}. Interpolating the \(p=10/3\) estimate with
the \(L^2\) identity gives every gain below \(1/4\) at \(p=3\).
Some three-dimensional metrics require
\(p\ge10/3\) for gains arbitrarily close to \(1/p\)
\cite{MinicozziSogge1997}, whereas the Euclidean exponent is \(3\).


\subsection{An estimate on one frequency annulus}\label{subsec:annular}

The proof is built around a single estimate at \(p=\pc\). It quantifies
the growth with frequency that the arbitrarily small Sobolev loss in
Theorem~\ref{thm:local-smoothing} allows. Our Fourier convention is
\(\widehat f(\xi)=\int_{\RR^n}e^{-\mathrm iX\cdot\xi}f(X)\dd X\),
so \(U(T)\) has phase \(X\cdot\xi-T|\xi|\) and inverse Fourier factor
\((2\pi)^{-n}\).

\begin{theorem}[Critical annular estimate]\label{thm:main}
Let \(n\ge2\), \(p=\pc\), and \(\varphi\in C_c^\infty(\RR^n\setminus\{0\})\).
For \(N\ge2\) set \(P_N=\varphi(D/N)\). Then
\begin{equation}\label{eq:main}
 \|U(T)P_N f\|_{L^p(\RR^n\times I)}
       \le C_{n,\varphi}[\log(2+N)]^{2/p}\|f\|_{L^p(\RR^n)} .
\end{equation}
\end{theorem}

Interpolation and summation over annuli yield
Theorem~\ref{thm:local-smoothing} (Section~\ref{sec:transfer}).
For unrestricted data on a full annulus the optimal logarithmic power in \eqref{eq:main}
lies between \(1/\pc\) and \(2/\pc\). The lower bound comes from the
focusing example: every dyadic distance from the focus between one
wavelength and unit distance contributes comparable mass to the
\(\pc\)-th power of the spacetime norm (Proposition~\ref{prop:sharpness}).
With any positive angular regularity, the power \(1/\pc\) suffices;
Appendix~\ref{sec:angular} proves a sharper result with a logarithmic
angular weight. For unrestricted data the lower bound persists even
for the weak spacetime norm (Proposition~\ref{prop:weak-obstruction}).

\subsection{A conserved current and two circles}\label{sec:idea}

Consider a finite sum of plane waves with $|\xi|\asymp N$ whose
directions lie in one small cap. We need to take a square \emph{after} summing these waves:
summing the squares would discard the interference that produces a
focus. An imaginary translation along a light ray filters the sum by
direction. Translation by $-\mathrm i\ell(1,\omega)$, with $|\omega|=1$ and
$\ell>0$, multiplies a mode of frequency $\xi$ by
$\exp[-\ell|\xi|\,|\omega-\xi/|\xi||^2/2]$. Its Gaussian width
in direction is therefore of order $(N\ell)^{-1/2}$ for $N\ell\gg1$.
The symmetric first-order wave system provides these light rays as
its quadratic fluxes. The sum of two such fluxes gives the imaginary
displacement. Two flows built from the continued wave then give a
positive semidefinite tensor. Each of its real columns has zero divergence,
and its time component still detects the original wave on the chosen cap
(Section~\ref{sec:current}).

Here is the estimate that the construction must deliver. Set $h=N^{-1}$,
let $a=Kh$ for a fixed large $K$, and write $p=\pc$. In addition to the
physical time $T$, the construction has two auxiliary time coordinates
$\tau\in\RR^2$. Write $H(T,s)$ for the circular mean at $|\tau|=s$
of the nonnegative energy $\mathcal F(T,\tau)$ constructed below. For $h$ small and nonzero
cap data on a fixed spatial torus normalized by $\|f\|_2=1$, we prove
\begin{align*}
 \int_1^2\|U(T)f\|_p^p\dd T
 &\ \lesssim\ a^{-2}\!\int_1^2 H(T,2a)\dd T\\
 &\ \lesssim\ [\log(2/h)]^2 a^{-2} H(0,4a)
 \ \lesssim\ [\log(2/h)]^2\|f\|_p^p .
\end{align*}
The first and last inequalities compare the energy with the wave. The
middle one uses only a scalar differential inequality.

To obtain the energy, fix $(T,\tau)$ and let the remaining real and
auxiliary spatial variables range over a bounded fiber with periodic
real directions. Write $\mu>0$ for the time component of the regularized
current. There is a complex scalar $\Phi$ constant along the current.
With $\lambda>0$ large, take on each fiber the solution with positive
definite spatial Hessian $H_{\rm fib}(\lambda|\Phi|^2+u)$ of
\[
   \det H_{\mathrm{fib}}(\lambda|\Phi|^2+u)=\lambda\mu,
   \qquad u|_{\partial}=0 .
\]
Here $H_{\rm fib}$ is an $n$-by-$n$ quaternionic Hermitian Hessian; its
determinant on positive matrices is the product of the eigenvalues.
The background costs nothing in the conservation law: since the
current is tangent to the levels of $\Phi$, its pairing with
$H(|\Phi|^2)$ vanishes. In the spatial determinant, however, it
supplies one direction of curvature. The Hessian of $\lambda|\Phi|^2$ has just one nonzero eigenvalue,
of size $\lambda$. Adding $\alpha I$ to this background changes its
determinant to $(\beta+\alpha)\alpha^{n-1}$, where $\beta\asymp\lambda$.
Thus density $\mu\ge m$ on a small ball forces the depth of $u$ at
its center to be at least a constant times the squared radius times
$m^{1/(n-1)}$, in the range of the construction. The integral
$\int(-u)\mu$ detects the power $m^{n/(n-1)}$. Since $\mu$ recovers
a square of the wave, the resulting exponent is $2n/(n-1)$.

A correction to $\int(-u)\mu$ makes the full determinant cancel when
one differentiates the energy in $(T,\tau)$, while leaving the corrected
energy comparable to that integral. The correction uses the integral
of the spatial determinant along the segment of potentials from
$\lambda|\Phi|^2$ to $\lambda|\Phi|^2+u$, as in the classical
Aubin--Mabuchi energy \cite{Aubin1984,Mabuchi1986}. The two integrations by parts use conservation of the
wave current and of the Hessian cofactor (Sections~\ref{sec:calculus}
and~\ref{sec:energy}). The cofactor contributes no boundary term:
its normal--normal coefficient vanishes.
Only the wave flux through the side of the fiber can give a positive
contribution to the Laplacian of the energy. A comparison function made
from a spatial average of the squared wave bounds $\mathcal F$ on the
initial circle and the normal derivative of $u$ at the side of the fiber.
On the part of the boundary that carries wave flux, the Gaussian decay
is stronger than the cost of recovering the wave in the interior.
The resulting scalar inequality is
\[
 \Delta_{T,\tau}\mathcal F\le C s^{-2}e^{-cs/h}\mathcal F,
       \qquad 2a\le s=|\tau|\le3
       \quad\text{(Section~\ref{sec:barrier}).}
\]

It remains to compare the circles. The scalar inequality carries a
lower bound at the observed times $1\le T\le2$ from $|\tau|=2a$
to $|\tau|=1$. At that fixed radius the mean inequality carries it to
initial time; a second annulus comparison then takes it back to
$|\tau|=4a$. A radial harmonic function in the plane is affine in the
logarithm of its radius. Each passage thus loses an inverse logarithm.
Both losses can occur for a scalar supersolution even with zero
potential (Remark~\ref{rem:cylinder-power}). Taking a $p$th root and
summing a fixed number of caps gives the claimed logarithmic bound on
the torus. Section~\ref{sec:transfer} removes the spatial period to prove
Theorem~\ref{thm:main}, then sums the frequencies for the Sobolev estimate.

\paragraph{Related constructions.}
The bilinear flux of a symmetric hyperbolic system is classical
\cite{Friedrichs1954}. Analytic continuation in an imaginary spacetime
direction has also been used to study waves \cite{KaiserStreater1992}.
For positive-energy Klein--Gordon waves, Kaiser found a vector
conserved in physical spacetime at each fixed imaginary displacement.
Integration over displacements gives a current with nonnegative time density
\cite[Section~4.5, Equations~(36)--(40)]{Kaiser1990}. Here the continued wave gives a positive
symmetric tensor whose columns have zero divergence in the combined
physical and auxiliary variables.

Serre's compensated integrability estimates pair a positive
divergence-free tensor with the Hessian of a real Monge--Ampère
potential to bound a power of its determinant
\cite[Section~5.1]{Serre2018}. The usual stress tensor of a wave need
not be positive in \(n\ge2\) spatial dimensions
\cite[Section~4.1]{Serre2019}. Our lifted tensor is positive
semidefinite, but before averaging it is a sum of four real squares
and its determinant vanishes. We use its regularized time component as the prescribed
density in a quaternionic determinant equation on each spatial fiber
(Section~\ref{sec:energy}).

The main proof ends in Section~\ref{sec:transfer}; Appendix~\ref{sec:dirichlet}
supplies the Dirichlet theorem used in the construction.
Sections~\ref{sec:kakeya}--\ref{sec:parabolic} give the applications.
Appendix~\ref{sec:planar} replaces the determinant equation by a linear
Dirichlet problem on a disk in two dimensions. Appendix~\ref{sec:pencil}
extends the current construction to a simple extremal eigenvalue of a
real symmetric \(n\)-by-\(n\) pencil in \(n\) spatial variables whose
transverse Hessian is definite, and gives an example with an irreducible
cubic characteristic cone.
